\input{../tex-inputs/latex-leseansicht-vorspann}

\section[Olga und Arthur Schnitzler an Gustav Schwarzkopf, 28. .5. 1916]{L04405 Olga und Arthur Schnitzler an Gustav Schwarzkopf, 28. .5. 1916}
\nopagebreak\mylabel{L04405v}
\rehead{ }\normalsize\beginnumbering\briefempfaengerindex{Schwarzkopf, Gustav@\textsc{Schwarzkopf, Gustav}!zzzSchnitzler, Arthur@\emph{von Arthur Schnitzler}!1916-05-281@{28. .5. 1916}|(be}\briefempfaengerindex{Schwarzkopf, Gustav@\textsc{Schwarzkopf, Gustav}!zzzSchnitzler, Olga@\emph{von Olga Schnitzler}!1916-05-281@{28. .5. 1916}|(be}
\toendnotes[C]{\smallbreak\pagebreak[2]}
\correspDesc{Versand  durch Olga Schnitzler, Arthur Schnitzler am 28. .5. 1916 in Bad Ischl
\newline{}Übermittlung  am 29. .5. 1916 in Bad Ischl
\newline{}Erhalt  durch Gustav Schwarzkopf im Zeitraum [29. 5. 1916 – 2. 6. 1916?] in Wien}\toendnotes[C]{\smallbreak}
\Standort{DLA, A:Schnitzler, HS.1985.1.1897.}
\physDesc{Bildpostkarte, 456 Zeichen
\newline{}Handschrift Olga Schnitzler: Bleistift, lateinische Kurrent
\newline{}Handschrift Arthur Schnitzler: Bleistift, deutsche Kurrent}\toendnotes[C]{\smallbreak}\pstart{}{\pb}Herrn Gustav Schwarzkopf\pend{}\pstart{}Wien I\oindex{I., Innere Stadt@\textbf{I., Innere Stadt}|pw}.\pend{}\pstart{}Tiefer Graben 17\oindex{Wien@\textbf{Wien}!I., Innere Stadt@\textbf{I., Innere Stadt}!Tiefer Graben 17@\textbf{Tiefer Graben 17}|pw}.\pend{}{\bigskip}
\pstart
           \centering{}{\pb}\textcolor{gray}{\textbf{Bad-Ischl\oindex{Bad Ischl@\textbf{Bad Ischl}|pw}.}}\textcolor{gray}{\textbf{Jagdgruppe\pwindex{Lauscher@\emph{Lauscher}|pwv} bei der kaiserl.
                  Villa\oindex{Kaiservilla@\textbf{Kaiservilla}|pw}.}}\pend
           \vspace{1em}
\pstart
           {\pb}Herzliche
               Grüsse{\\[\baselineskip]}von Ihrer \spacefill\mbox{Olga S.}\pend
           \leftskip=0em{}\selectlanguage{ngerman}\vspace{1em}
\pstart
           \raggedleft{}{\pb}{[}hs. A. Schnitzler:{]} 28. 5. 16\pend
           \vspace{0.5em}
\pstart
           lieber Gustav, wir wohnen in den zwei{ }ſchönen Zimmern \substVorne{}\textsuperscript{bei}\substDazwischen{}im\substHinten{}{ }Athen\oindex{Hotel garni Athen@\textbf{Hotel garni Athen}|pw} – halten es übrigens als einzige  beſetzt,
               ohne andre Berührungspunkte mit der Entente. Im übrigen regnet es mit beka{\geminationn}ter Salzka{\geminationm}ergüte\oindex{Salzkammergut@\textbf{Salzkammergut}|pw} – wir gehn trotzdem vergnügt{ }ſpazieren – es iſt
               menſchenleer – auch bei So{\geminationn}enſchein\oindex{Restaurant Sonnenschein@\textbf{Restaurant Sonnenschein}|pw} und Zauner\oindex{Konditorei Zauner@\textbf{Konditorei Zauner}|pw}, die auf alter Höhe beharren.
                  \label{K_L04405-1v}\edtext{Dienſtag fahren wir nach Altauſſee\oindex{Altaussee@\textbf{Altaussee}|pw}}{\lemma{\textnormal{\emph{Dienstag … Altaussee}}}\Cendnote{\textnormal{Siehe A. S.: \emph{Wiener Schnitzler}, 30. 5. 1916.}}}\label{K_L04405-1}.\pend
           \pstart Herzlichſte Gr \spacefill\mbox{A.}\pend{}\selectlanguage{ngerman}\endnumbering\briefempfaengerindex{Schwarzkopf, Gustav@\textsc{Schwarzkopf, Gustav}!zzzSchnitzler, Arthur@\emph{von Arthur Schnitzler}!1916-05-281@{28. .5. 1916}|)be}\briefempfaengerindex{Schwarzkopf, Gustav@\textsc{Schwarzkopf, Gustav}!zzzSchnitzler, Olga@\emph{von Olga Schnitzler}!1916-05-281@{28. .5. 1916}|)be}\mylabel{L04405h}
\begin{anhang}
\end{anhang}\newcommand{\dateiname}{L04405}\newcommand{\titel}{Olga und Arthur Schnitzler an Gustav Schwarzkopf, 28. .5. 1916}\newcommand{\editorInnen}{Herausgegeben von Jahnke, SelmaMüller, Martin Anton}\input{../tex-inputs/latex-leseansicht-abspann}
